% Final paper of the TCL research program (Phases 1-4a), TMLR format.
% Compile with the official TMLR style files (tmlr.sty, tmlr.bst, fancyhdr.sty from
% https://github.com/JmlrOrg/tmlr-style-file) next to this file, e.g. on Overleaf.
% Anonymous TMLR submission: \usepackage{tmlr} (and replace the author block and the
%   repository links with anonymized ones).
% arXiv preprint: \usepackage[preprint]{tmlr}.
\documentclass[10pt]{article}
\usepackage[preprint]{tmlr}

\usepackage{hyperref}
\usepackage{url}
\usepackage{booktabs}
\usepackage{amsmath}
\usepackage{multirow}

\newcommand{\Ltrust}{\mathcal{L}_{\text{trust}}}
\newcommand{\LLM}{\mathcal{L}_{\text{LM}}}

\title{Token Confidence Is Hard to Beat:\\A Preregistered Study of Learned Trust Heads\\for Language Model Question Answering}

\author{\name Awab Mohamed \\
  \addr Independent researcher, Qatar\\
  \url{https://github.com/awabmoha/truth-calibration-layer}}

\def\month{MM}
\def\year{2026}
\def\openreview{\url{https://openreview.net/forum?id=XXXX}}

\begin{document}
\maketitle

\begin{abstract}
Trained probes on the hidden states of large language models (LLMs) are widely reported to detect incorrect answers better than the model's own token probabilities. Such comparisons are rarely made on a dataset the probe was not trained on, and rarely under a decision rule fixed in advance. In this work, we test learned trust heads against the answer-span token probability on held-out short-answer question answering (QA) datasets, in a preregistered program of four phases with three models (Qwen2.5-0.5B, Phi-3.5-mini, and Qwen2.5-7B). Every protocol was locked before its data existed, and every analysis was run once. We first show that frozen probes do not transfer across datasets, and that training the head jointly with the backbone does not make it transfer better, in two independent runs. A head trained on a mixture of datasets, and given the token probability as an input, beats that probability on 1 of 6 held-out pairs, and on none of 3 at 7B. Two further trust dimensions, self-consistency and passage grounding, are learnable from a single forward pass but add little to token confidence; the grounding head mainly detects answers given from irrelevant passages, which the 7B model refuses by itself. Finally, exploratory analyses show that the outcome of such comparisons depends on the baseline: against the first-token probability the learned head wins on 4 of 9 held-out pairs, against the best of three free token-probability baselines on 1 of 9. These results do not establish a new phenomenon; they subject claims made under weaker conditions to a strict, fully reported test.
\end{abstract}

\section{Introduction}

Language models produce wrong answers as fluently as right ones.\footnote{The author used an AI assistant (Claude, Anthropic) to help write the code and to draft and edit the text. The author designed the study, made all protocol and analysis decisions, and reviewed every result and the final text, and is responsible for the ideas, claims, and results.} An estimate of how likely each answer is to be correct allows a system to show, flag, or withhold it. The model's own token probabilities provide such an estimate at no cost. A large body of work trains estimators that are reported to do better: probes and heads on hidden states \citep{azaria-mitchell-2023-internal,marks2024geometry,orgad2025llms,liu2024universal}, models fine-tuned to estimate their own uncertainty \citep{kapoor2024large}, and one-pass predictors of sampling-based uncertainty \citep{kossen2024semantic,vashurin2025cocoa}.

The evidence on whether these estimators generalize is mixed. \citet{kapoor2024large} find that fine-tuned estimators transfer to new question types, while \citet{orgad2025llms} find that probe transfer rarely exceeds what logits alone achieve, and LM-Polygraph \citep{vashurin2025benchmarking} finds sequence probability difficult to surpass on short outputs. These studies differ in their baselines, their training data, and whether test data come from a dataset the estimator was trained on, which makes the disagreement hard to resolve.

In this work, we test the strongest designs we could build against the answer-span token probability, on datasets held out from training, with every hypothesis, metric, and pass rule fixed before the corresponding data existed. The tests were organized as a research program of four phases (Section~\ref{sec:program}). We ask whether a learned correctness head beats token confidence on a new dataset (Section~\ref{sec:correctness}), whether training the head jointly with the backbone improves its transfer (Section~\ref{sec:joint}), and whether two further trust dimensions, self-consistency and grounding, add to token confidence (Sections~\ref{sec:consistency} and~\ref{sec:grounding}).

We find that, on short-answer QA from 0.5B to 7B parameters, learned hidden-state signals add little or nothing beyond the answer-span probability on held-out data. Frozen probes outperform it in distribution only where token confidence is poor, and do not transfer. Joint training improves the head in distribution but not on a new dataset. Mixture training and access to the token probability do not close the gap, and at 7B the head is slightly worse than its own input. Self-consistency predicted from one pass is the most robust addition, matching 10-sample semantic entropy on several pairs, but its gain in selective prediction passes the preregistered rule at no model size. The grounding head mostly detects irrelevant passages, a case the 7B model handles by abstaining, and its verdict depends on how the grounding label is constructed. Exploratory analyses (Section~\ref{sec:exploratory}) show that the choice among free token-probability baselines can reverse the comparison, which may explain part of the disagreement in prior work (Section~\ref{sec:discussion}).

\subsection{Contributions}

None of these findings is a new phenomenon (Section~\ref{sec:related}); our contributions concern the strictness and completeness of the test.
\begin{enumerate}
\item We test learned correctness heads, including a head given the token probability as an input, against the answer-span probability on held-out datasets at three model sizes.
\item We compare joint, trust-only, LM-only, and frozen training in a controlled four-arm design, run twice.
\item We test one-pass self-consistency and grounding heads as additions to token confidence across datasets, with label audits before every analysis.
\item We provide diagnostics that bear on conflicting published results: the effect of the baseline choice, the dependence of gains on weak token confidence, and a labeling choice that reverses a preregistered verdict.
\item We release every protocol, amendment, audit, record, and analysis, including all failed tests.
\end{enumerate}

\section{Related Work}
\label{sec:related}

\paragraph{Probes and their transfer.} Classifiers on hidden states detect false statements and incorrect answers \citep{azaria-mitchell-2023-internal,burns2023discovering,marks2024geometry}, and models can be trained to predict whether they know an answer \citep{kadavath2022language}. Transfer to new tasks is weaker: such predictions only partly generalize \citep{kadavath2022language}, probes are sensitive to their training domain \citep{chwang2024androids}, and after subtracting the strongest logit baseline, probe generalization rarely exceeds what logits achieve \citep{orgad2025llms}. Probe advantages over log-probability baselines concentrate in in-domain and within-family transfer \citep{duan2026probes}. Training on many datasets improves transfer \citep{liu2024universal}, and estimators combining logit and hidden-state features are reported to transfer \citep{liu2024supervised}; trained probes were the most reliable estimators in a comparison by \citet{mahaut2024factual}. Our study differs in using a span-level token-probability baseline on a dataset never used for training or design, under a rule fixed in advance.

\paragraph{Training the model for uncertainty.} \citet{kapoor2024large} fine-tune models with LoRA on about 20k graded examples from 16 benchmarks and find that training through the model's features is necessary for uncertainty estimates that generalize, while a frozen probe is not sufficient. Logit-based uncertainty generalizes under shift about as well as trained verbal confidence \citep{lin2022teaching}, supervised uncertainty heads work best in domain \citep{shelmanov2025head}, and fine-tuning can distort pretrained features and underperform linear probing out of distribution \citep{kumar2022finetuning}. RAGognizer \citep{ridder2026ragognizer} trains a detection head jointly with LoRA adapters. We isolate the joint objective by comparing it with trust-only, LM-only, and frozen training on identical data.

\paragraph{Sampling-based uncertainty and its amortization.} Semantic entropy clusters sampled answers by meaning \citep{kuhn2023semantic,farquhar2024detecting}. Semantic Entropy Probes predict it from one forward pass and generalize better than accuracy probes, but do not match sampling \citep{kossen2024semantic}. CoCoA combines token-probability confidence with agreement across samples; its Light variant predicts that agreement from mid-layer hidden states and, combined with sequence probability, outperforms semantic entropy \citep{vashurin2025cocoa}. Our self-consistency head follows the same idea; we evaluate it on held-out datasets against a preregistered selective-prediction criterion.

\paragraph{Grounding and irrelevant context.} Contextual hallucinations can be detected from attention patterns \citep{chuang2024lookback} and model internals \citep{sun2025redeep}, but representation-based detectors often fail out of distribution, partly because of spurious correlations in their training data \citep{dubanowska2025representation}. Irrelevant passages change model behavior \citep{cuconasu2024power,yoran2024robust}, and strong models often answer when the context is insufficient \citep{joren2025sufficient}. Converting a question and answer into an NLI hypothesis is error-prone \citep{demszky2018qa2d,chen-etal-2021-nli}. We separate a grounding head's gains by passage type and measure the effect of the hypothesis construction on the verdict.

\paragraph{Simple baselines and preregistration.} In the LM-Polygraph benchmark, maximum sequence probability is a strong and robust baseline that is often difficult to surpass on short outputs \citep{vashurin2025benchmarking}. Preregistration has been argued for in NLP \citep{vanmiltenburg-etal-2021-preregistering,sogaard-etal-2023-two}; a closely related precedent is a single-model preregistered negative result on distilling self-consistency into verbal confidence \citep{cacioli2026distilling}. We apply preregistration across a multi-phase program and report every outcome.

\section{Preregistered Program}
\label{sec:program}

The study was run as a sequence of phases. Each phase had a protocol specifying hypotheses, arms, metrics, a size rule, and a pass rule, committed before any data for the phase was generated or any test metric computed. Each analysis was run once. Label audits on random samples were completed before any metric was read. Changes after locking were recorded as dated amendments, none of which altered a decision rule after a test metric existed (Appendix~\ref{app:amendments}). Table~\ref{tab:program} lists all ten decisions. Two planned phases, a non-QA task family with joint training at 7B and claim-level trust in long-form text, were not run.

\begin{table}[t]
\centering
\small
\begin{tabular}{@{}l p{2.6cm} p{6.2cm} p{3.9cm}@{}}
\toprule
Phase & Hypothesis & Test & Result \\
\midrule
1 & H1-R & Frozen probe beats token confidence, in distribution & Passes, at the threshold (4 of 6) \\
1 & H2-R & Frozen probe transfers across datasets & Fails (2 of 6 directions) \\
2 & H3-T, H3-J & Joint training transfers better than frozen / trust-only & Fail \\
2b & H3-T, H3-J & The one allowed rerun, after a bug fix & Fail again \\
3 & H4-T & Mixture-trained head beats token confidence on a held-out set & Fails (1 of 6) \\
3 & H4 (consistency) & One-pass self-consistency adds to token confidence & Fails by one pair (3 of 6) \\
3 & H4 (grounding) & One-pass grounding adds to token confidence & Fails (2 of 7) \\
4a & H4-T at 7B & As Phase 3, Qwen2.5-7B & Fails (0 of 3) \\
4a & H4 (consistency) at 7B & As Phase 3 & Fails (2 of 3, without PopQA) \\
4a & H4 (grounding) at 7B & As Phase 3 & Not testable \\
\bottomrule
\end{tabular}
\caption{All preregistered decisions, in the order they were made. Counts give the pairs (model and held-out dataset) or transfer directions that met the pass criterion, out of those eligible under the size rule.}
\label{tab:program}
\end{table}

\section{Setup}
\label{sec:setup}

\paragraph{Models.} We used Qwen2.5-0.5B-Instruct and Qwen2.5-7B-Instruct \citep{qwen2024qwen25}, Phi-3.5-mini-instruct \citep{abdin2024phi3}, and, in Phase 1 only, Gemma-2B-it \citep{gemma2024}. Each model answered with its chat template and greedy decoding; the 7B model ran in float16 across two T4 GPUs.

\paragraph{Datasets.} Closed-book QA used TriviaQA \citep{joshi-etal-2017-triviaqa}, SQuAD \citep{rajpurkar-etal-2016-squad} with its passage and split by Wikipedia article, and NQ-Open \citep{kwiatkowski-etal-2019-natural}, each with 2,500 training, 500 validation, and 1,500 test questions. PopQA \citep{mallen-etal-2023-trust} served as a dataset never used for training or design (1,500 test questions; 3,000 for Qwen2.5-0.5B, extended under the preregistered size rule). Grounding used the MRQA \citep{fisch-etal-2019-mrqa} versions of SQuAD, Natural Questions, and HotpotQA \citep{yang-etal-2018-hotpotqa} (2,500/500/1,500 questions), with RelationExtraction (1,500 test questions) held out. In every split, one third of the questions were shown an unrelated passage from the same dataset (a distractor), and the prompt permitted the answer ``unanswerable''.

\paragraph{Labels.} An answer was labeled correct if it matched an accepted answer under a fixed normalization rule, which was validated on audited answers before Phase 2 and audited again in every phase (Appendix~\ref{app:audits}). The self-consistency label was the share of 10 answers sampled at temperature 1 that an NLI model \citep{he2021deberta} judged equivalent to the greedy answer. The grounding label was NLI entailment of the answer by the passage the model saw, at a threshold of 0.5.

\paragraph{Baseline.} The baseline R is the answer-span probability: the geometric mean of the probabilities of the answer tokens, mapped to a probability by Platt scaling \citep{platt1999probabilistic} on validation data. We chose it because it needs no training and no sampling. Platt scaling is monotone and therefore leaves ranking metrics unchanged; what distinguishes this baseline from many in prior work is the aggregation over the answer span rather than the calibration.

\paragraph{Trust heads.} Phase 1 fitted a logistic probe on final-layer hidden states averaged over the answer span. Phase 2 used a two-layer MLP on the 50\%-depth answer-mean state, in four arms trained on identical data: (A) a frozen backbone; (B) LoRA \citep{hu2022lora} trained on $\Ltrust$; (C) LoRA trained on $\LLM + \lambda\Ltrust$; and (D) LoRA trained on $\LLM$, followed by a head. Phases 3 and 4a used a single MLP design (LayerNorm, 256 hidden units, GELU, weight decay 1, equal weight per training dataset), selected once on validation splits and trained leave-one-dataset-out with 5 seeds. M denotes the head on hidden states and M+R the head that also receives the logit of R. For the additional dimensions, a one-pass head predicts self-consistency (C-hat) or grounding (G-hat), and a logistic combiner fitted on validation data merges it with R (R+C-hat, R+G-hat). Abstentions were excluded from every grounding arm.

\paragraph{Metrics.} We measure correctness ranking by AUROC and selective-prediction risk by the area under the risk--coverage curve (AURC; \citealp{geifman2017selective}). For two scores $a$ and $b$, $d(a,b)$ denotes the difference in AUROC averaged over seeds; for an added dimension $X$, the AURC gain is $\mathrm{AURC}(\mathrm{R}) - \mathrm{AURC}(\mathrm{R}{+}X)$, so that positive values favor the combination. Confidence intervals are 95\% intervals from 1,000 paired bootstrap resamples of test items.

\paragraph{Decision rules.} H4-T passes on a pair if the confidence interval of $d(\text{M+R}, \text{R})$ lies above zero. A trust dimension passes on a pair if its head's AUROC for its own label is at least 0.70, its Spearman correlation with R is below 0.7, and the confidence interval of its AURC gain lies above zero. A test passes if a majority of its eligible pairs pass, including the pair for the new dataset (PopQA, or RelationExtraction for grounding). A pair is eligible only if each class has at least 150 test answers.

\section{Results}

\subsection{Learned correctness heads do not beat token confidence on held-out data}
\label{sec:correctness}

\paragraph{In distribution and across one dataset shift.} In Phase 1, a frozen probe outperforms R on 4 of 6 pairings, exactly the preregistered threshold. The gains are +0.064 and +0.034 AUROC for Qwen on TriviaQA and Phi on NQ-Open, and +0.306 and +0.250 for Gemma, whose token confidence is close to chance (AUROC 0.527 and 0.620); on Qwen SQuAD the probe is significantly worse ($-$0.054). Transferred to the other dataset, the probes underperform R in 3 of 4 Qwen and Phi directions (e.g., 0.602 vs.\ 0.772 for Qwen from TriviaQA to SQuAD) and match it in the fourth. Only Gemma's probes transfer (+0.21 and +0.25), consistent with the weakness of its token confidence.

\paragraph{With mixture training and access to the token probability.} Table~\ref{tab:h4t} shows that training on a mixture of datasets and supplying R as an input does not close the gap. The head beats R on 1 of 6 pairs in Phase 3 and on none of 3 at 7B. Mixture training does improve on single-dataset heads on every Phase 3 pair, which indicates that the failure is not due to undertrained heads. At 7B, adding hidden states to R lowers AUROC on all three pairs, with every seed below zero on TriviaQA and PopQA, although no individual difference is significant. The head thus fits patterns in the training datasets that do not carry over to the held-out one.

\begin{table}[t]
\centering
\small
\begin{tabular}{@{}llrrrrr@{}}
\toprule
Model & Held out & R & M & M+R & $d(\text{M+R}, \text{R})$ [95\% CI] & SE \\
\midrule
\multirow{3}{*}{Qwen 0.5B} & TriviaQA & 0.712 & 0.740 & 0.744 & +0.032 [$-$0.010, +0.077] & 0.773 \\
 & SQuAD & 0.762 & 0.624 & 0.667 & $-$0.096 [$-$0.131, $-$0.061] & 0.775 \\
 & PopQA & 0.828 & 0.909 & 0.907 & \textbf{+0.079} [+0.052, +0.105] & 0.879 \\
\midrule
\multirow{3}{*}{Phi 3.8B} & TriviaQA & 0.836 & 0.795 & 0.795 & $-$0.041 [$-$0.064, $-$0.017] & 0.842 \\
 & NQ-Open & 0.765 & 0.759 & 0.760 & $-$0.006 [$-$0.036, +0.024] & 0.762 \\
 & PopQA & 0.864 & 0.815 & 0.809 & $-$0.055 [$-$0.080, $-$0.029] & 0.881 \\
\midrule
\multirow{3}{*}{Qwen 7B} & TriviaQA & 0.854 & 0.832 & 0.833 & $-$0.021 [$-$0.042, +0.000] & 0.864 \\
 & NQ-Open & 0.742 & 0.734 & 0.738 & $-$0.004 [$-$0.035, +0.027] & 0.747 \\
 & PopQA & 0.865 & 0.834 & 0.844 & $-$0.022 [$-$0.045, +0.005] & 0.884 \\
\bottomrule
\end{tabular}
\caption{Learned correctness heads against the answer-span probability on held-out datasets (H4-T). Values are test AUROC (higher is better), averaged over 5 seeds. R: answer-span probability; M: mixture-trained head; M+R: head that also receives R; SE: 10-sample semantic entropy, shown for reference and not gated. Bold marks the only difference whose confidence interval lies above zero. Pairs failing the size rule (Qwen 0.5B on NQ-Open, SQuAD at 7B) are omitted.}
\label{tab:h4t}
\end{table}

\subsection{Joint training improves the head in distribution but not on a new dataset}
\label{sec:joint}

Table~\ref{tab:joint} reports transfer AUROC for the four Phase 2 arms. Joint training (C) improves on the frozen head (A) from TriviaQA to SQuAD by 0.022 AUROC but is significantly worse from SQuAD to TriviaQA, by 0.038. LM-only fine-tuning (D) improves on A by 0.031 in the first direction, matching the gain of joint training, which suggests that the improvement comes from language-model fine-tuning rather than from the joint objective. After the analysis, we found that the language-modeling target omitted the end-of-turn token, which made fine-tuned models produce long answers. The one rerun permitted by the protocol (Phase 2b) restored answer length but not the transfer result (+0.034 and $-$0.038). In both runs, every trained arm transfers below the untrained token probability and semantic entropy.

\begin{table}[t]
\centering
\small
\begin{tabular}{@{}lrrrr@{}}
\toprule
 & \multicolumn{2}{c}{Phase 2} & \multicolumn{2}{c}{Phase 2b} \\
Arm & TQA$\to$SQuAD & SQuAD$\to$TQA & TQA$\to$SQuAD & SQuAD$\to$TQA \\
\midrule
A: frozen head & 0.557 & 0.645 & 0.557 & 0.645 \\
B: trust-only LoRA & 0.572 & 0.629 & 0.572 & 0.629 \\
C: joint LoRA & 0.578 & 0.606 & 0.591 & 0.606 \\
D: LM-only LoRA, then head & 0.588 & 0.643 & 0.598 & 0.647 \\
\midrule
Token probability (untrained) & 0.762 & 0.712 & 0.762 & 0.712 \\
Semantic entropy (untrained) & 0.775 & 0.773 & 0.775 & 0.773 \\
\bottomrule
\end{tabular}
\caption{Joint training against frozen, trust-only, and LM-only training (Qwen2.5-0.5B). Values are transfer AUROC (higher is better), averaged over 3 seeds, on 1,500 fixed answers from the dataset not used for training. Arms A and B were reused in Phase 2b; C and D were retrained after the end-of-turn fix. No trained arm reaches the untrained baselines.}
\label{tab:joint}
\end{table}

\subsection{Self-consistency is learnable from one pass but adds little}
\label{sec:consistency}

The consistency head predicts its label on every held-out pair (AUROC 0.73--0.95) without duplicating R (Spearman 0.26--0.64). Combined with R, the one-pass prediction matches or exceeds 10-sample semantic entropy on several pairs, for example 0.785 vs.\ 0.773 for Qwen 0.5B on TriviaQA and 0.772 vs.\ 0.747 for Qwen 7B on NQ-Open, in line with CoCoA Light \citep{vashurin2025cocoa}. Its gain in selective prediction, however, passes on 3 of 6 pairs at 0.5B and 3.8B and on 2 of 3 at 7B (Table~\ref{tab:dims}). At 7B the PopQA pair, required by the rule, misses by a margin of 0.0003 at the lower confidence bound. Within the Qwen family, the gain decreases from about 0.020 AURC at 0.5B to 0.007--0.009 at 7B on TriviaQA and PopQA.

\begin{table}[t]
\centering
\small
\begin{tabular}{@{}lllll@{}}
\toprule
 & \multicolumn{2}{c}{Self-consistency (R + C-hat)} & \multicolumn{2}{c}{Grounding (R + G-hat)} \\
Model & Held out & AURC gain [95\% CI] & Held out & AURC gain [95\% CI] \\
\midrule
\multirow{4}{*}{Qwen 0.5B} & TriviaQA & \textbf{+0.021} [+0.010, +0.034] & SQuAD & \textbf{+0.009} [+0.001, +0.017] \\
 & SQuAD & +0.003 [$-$0.001, +0.008] & NQ & \textbf{+0.014} [+0.004, +0.024] \\
 & PopQA & \textbf{+0.020} [+0.015, +0.026] & HotpotQA & +0.009 [$-$0.000, +0.019] \\
 & & & RelationExt. & +0.004 [$-$0.002, +0.010] \\
\midrule
\multirow{3}{*}{Phi 3.8B} & TriviaQA & $-$0.008 [$-$0.018, +0.001] & SQuAD & +0.004 [$-$0.001, +0.010] \\
 & NQ-Open & \textbf{+0.016} [+0.007, +0.026] & NQ & $-$0.011 [$-$0.026, +0.000] \\
 & PopQA & +0.003 [$-$0.006, +0.011] & HotpotQA & +0.001 [$-$0.011, +0.014] \\
\midrule
\multirow{3}{*}{Qwen 7B} & TriviaQA & \textbf{+0.009} [+0.001, +0.017] & NQ & $-$0.010 [$-$0.015, +0.000] \\
 & NQ-Open & \textbf{+0.020} [+0.011, +0.029] & HotpotQA & $-$0.012 [$-$0.020, $-$0.003] \\
 & PopQA & +0.007 [$-$0.000, +0.015] & & \\
\bottomrule
\end{tabular}
\caption{Gains from adding a one-pass self-consistency or grounding head to the answer-span probability. Values are AURC(R) minus AURC of the combination (higher is better), averaged over 5 seeds. Bold marks gains whose confidence interval lies above zero. The new datasets are PopQA for self-consistency and RelationExtraction for grounding. Pairs failing the size rule are omitted; at 7B only NQ and HotpotQA were eligible for grounding.}
\label{tab:dims}
\end{table}

\subsection{Grounding heads mostly detect irrelevant passages}
\label{sec:grounding}

The grounding head is learnable on 6 of 7 Phase 3 pairs (AUROC 0.70--0.84) and only weakly correlated with R (Spearman 0.11--0.61), but its AURC gain passes on 2 of 7. Separating answers by passage type explains the gap. Among answers given from the correct passage, the combination improves AUROC over R by at most 0.003; its pooled gain comes from ranking answers given from a distractor passage, which are almost all incorrect, below the rest. This also accounts for the small gains on Phi, which abstains on most distractors and leaves 28--122 distractor answers per held-out set, compared with 287--359 for Qwen 0.5B. The 7B model abstains on 1,996 of the 2,000 held-out distractor rows and is rarely wrong on the questions it answers (67 and 75 incorrect answers on SQuAD and RelationExtraction, below the size rule), so the grounding test is not testable at 7B; both eligible pairs fail.

\paragraph{Sensitivity to the grounding label.} The NLI grounding label agrees with the audit labels on 77 of 100 Phase 3 answers and 38 of 50 Phase 4a answers, and most disagreements are supported short answers labeled unsupported. Before any 7B data existed, a replacement was selected among three NLI candidates on the Phase 3 audit set: the same model with the hypothesis template ``The answer to the question `$q$' is $a$.'', which agrees on 79 of 100. Re-running the Phase 3 grounding analysis with this label, otherwise unchanged and outside the preregistered decision, yields 5 of 7 passing pairs instead of 2, including RelationExtraction. The recorded decision stands: the gain still comes mainly from distractor rows, and the replacement was selected on 100 of the same answers. The construction of the NLI hypothesis therefore determined the outcome of a preregistered test.

\subsection{Exploratory analyses}
\label{sec:exploratory}

The following analyses were conducted after the program closed, on saved scores only, and do not affect any decision.

\paragraph{Choice of baseline.} Two further token-probability baselines require no extra computation: the probability of the first answer token and the minimum token probability within the answer. Table~\ref{tab:baselines} shows that the minimum token probability exceeds the answer-span probability on 4 of the 9 H4-T pairs, by up to 0.068 AUROC. Against the best of the three baselines, the learned head wins on 1 of 9 pairs, by 0.012; the only Phase 3 H4-T win (+0.079 over R) therefore reflects mainly a weak baseline. Against the first-token probability, the head wins on 4 of 9 pairs, by up to 0.051. A study using that baseline would thus report more probe wins from the same data.

\begin{table}[t]
\centering
\small
\begin{tabular}{@{}lrrrrr@{}}
\toprule
Pair & Answer span & First token & Minimum token & Head M+R & Head $-$ best \\
\midrule
Qwen 0.5B / TriviaQA & 0.712 & 0.746 & 0.773 & 0.744 & $-$0.030 \\
Qwen 0.5B / SQuAD & 0.762 & 0.717 & 0.750 & 0.667 & $-$0.096 \\
Qwen 0.5B / PopQA & 0.828 & 0.884 & 0.896 & 0.907 & +0.012 \\
Phi 3.8B / TriviaQA & 0.836 & 0.807 & 0.834 & 0.795 & $-$0.041 \\
Phi 3.8B / NQ-Open & 0.765 & 0.708 & 0.749 & 0.760 & $-$0.006 \\
Phi 3.8B / PopQA & 0.864 & 0.874 & 0.881 & 0.809 & $-$0.072 \\
Qwen 7B / TriviaQA & 0.854 & 0.826 & 0.853 & 0.833 & $-$0.021 \\
Qwen 7B / NQ-Open & 0.742 & 0.707 & 0.735 & 0.738 & $-$0.004 \\
Qwen 7B / PopQA & 0.865 & 0.879 & 0.882 & 0.844 & $-$0.038 \\
\bottomrule
\end{tabular}
\caption{Three free token-probability baselines and the learned head on the H4-T held-out pairs (exploratory). Values are test AUROC (higher is better); the head is averaged over 5 seeds. The last column is the head minus the best of the three baselines. Selecting the best baseline after seeing the data favors the baselines; the preregistered baseline is the answer span.}
\label{tab:baselines}
\end{table}

\paragraph{Weak token confidence.} The largest learned gains occur where token confidence is poor: Gemma in Phase 1 (token AUROC 0.53--0.62; probe gains of 0.25 to 0.31, and the only probes that transfer) and Qwen 0.5B on PopQA. Across the 9 eligible pairs per head in Phases 3 and 4a, the Spearman correlation between the AUROC of R and the head's AUROC gain is $-$0.43 for the correctness head and $-$0.60 for self-consistency. Part of such a correlation is mechanical, because estimation noise in the baseline enters both quantities with opposite signs. After normalizing gains by headroom and estimating the two quantities on disjoint halves of the test items, the correlation persists for self-consistency ($-$0.48 to $-$0.53) and weakens for the correctness head ($-$0.27); with 9 pairs, neither confidence interval excludes zero. For grounding the sign is positive (+0.63), since its gains arise from distractor answers rather than weak confidence. We therefore treat the pattern as a hypothesis.

\paragraph{Power.} From the recorded confidence intervals, the minimum effect detectable with 80\% power is 0.030--0.045 AUROC for H4-T at 7B, 0.017--0.033 AUROC for the Phase 2 comparisons, and about 0.01 AURC for the self-consistency and grounding tests. The null results thus exclude gains of practical size for learned correctness heads, but not small ones; the self-consistency result at 7B on PopQA (+0.0069 against a detectable effect of 0.011) is close to the threshold.

\section{Discussion and Conclusion}
\label{sec:discussion}

\paragraph{Relation to positive results.} Our findings conflict with several published results, and the differences in setup offer plausible explanations without implying error on either side. \citet{kapoor2024large} trained on about 20k examples from 16 benchmarks with 7B and 13B models, without a language-modeling term, and compared against an uncalibrated perplexity baseline; we trained on one to three datasets of 2,500 questions at 0.5B and compared against the answer-span probability. Since data diversity also helped in our Phase 3, much larger and more varied training sets may be what allows fine-tuned estimators to transfer. Training on many datasets improves probe transfer \citep{liu2024universal,liu2024supervised}; our mixture-trained heads likewise improve on single-dataset heads but do not reach the span-level baseline. CoCoA Light and Semantic Entropy Probes \citep{vashurin2025cocoa,kossen2024semantic} agree with our self-consistency results in direction; under a held-out protocol and a preregistered criterion, the gain is real but small and decreases with scale within one model family. Finally, Table~\ref{tab:baselines} suggests that comparisons against first-token probability or perplexity favor learned probes more than comparisons against span-level or minimum-token probability, consistent with the strength of sequence probability on short outputs \citep{vashurin2025benchmarking}.

\paragraph{Limitations.} All results concern short-answer QA, where token probability is known to be strong; long-form generation, reasoning, and code, where it is weaker, were not tested. The largest model has 7B parameters, and the three sizes come from two families, so scale is confounded with family and the within-family comparison has two points. Each head reads a single layer with mean pooling, chosen once on validation data; other designs, such as attention features, multiple layers, or token-level heads, may perform better. The distractors are unrelated passages, the easier case; passages that are related but insufficient were not tested. The same NLI model defines both the self-consistency and the grounding labels, and the grounding label is noisy. Phase 2 used a single 0.5B model with three seeds. The baseline comparison and the weak-confidence pattern are exploratory.

\paragraph{Conclusion.} On held-out short-answer QA from 0.5B to 7B parameters, learned hidden-state trust signals do not outperform the model's own answer probability when the comparison is fixed in advance. Joint training improves heads in distribution but not on new data. Self-consistency predicted from one pass is the most robust addition and approaches sampling-based semantic entropy at a fraction of the cost, but adds little to token confidence. Grounding heads primarily detect irrelevant passages, which a larger model refuses on its own. The gains that remain appear where token confidence is weak, and whether a study observes them can depend on which token-probability baseline it uses. We recommend that evaluations of learned trust estimators on short answers report span-level and minimum-token probabilities, calibrated on held-out data, as baselines.

\section*{Reproducibility Statement}

All protocols, amendments, checkpoints, label audits, analysis scripts, score files, and model outputs are available at \url{https://github.com/awabmoha/truth-calibration-layer} (release v1.0.0). Hidden-state arrays are not included but can be regenerated with the provided Kaggle notebooks, which are pinned to specific commits. All analyses run on a CPU from the released files.

\bibliography{references}
\bibliographystyle{tmlr}

\appendix

\section{Design Changes and Amendments}
\label{app:amendments}

Two design changes were made when protocols were locked. Phase 2b is a separate protocol for the single rerun permitted by the Phase 2 protocol; it adds the end-of-turn token to the language-modeling target. The Phase 3 grounding protocol replaced the originally planned SQuAD 2.0 unanswerable questions with MRQA datasets and distractor passages, so that the test could hold out a dataset. Two amendments were logged after locking, each before the affected metrics existed. In Phase 2, the micro-batch size for SQuAD was reduced after an out-of-memory failure, leaving the effective batch size unchanged. In Phase 4a, the grounding audit took all 4 answered distractor rows and 46 rows with the correct passage, instead of 25 of each, because the 7B model abstained on almost every distractor. The full texts are part of each protocol.

\section{Label Audits}
\label{app:audits}

\begin{table}[h]
\centering
\small
\begin{tabular}{@{}llr@{}}
\toprule
Phase & Audit & Agreement with audit labels \\
\midrule
3 & PopQA correctness, 100 answers per model & 99 of 100 (each model) \\
3 & Self-consistency count, 50 answers per model & binary label 92 of 100 \\
3 & Grounding label, 50 answers per model & 77 of 100 \\
4a & PopQA correctness, 100 answers & 99 of 100 \\
4a & Self-consistency count, 50 answers & exact count 43 of 50; binary label 47 of 50 \\
4a & Grounding label, 50 answers & 38 of 50 (11 missed support, 1 false support) \\
4a & Correctness on grounding answers, 50 answers & 47 of 50 \\
\bottomrule
\end{tabular}
\caption{Label audits on random samples, each completed before any test metric of its phase was read. The binary self-consistency label counts an answer as consistent when at least 5 of 10 samples agree with it.}
\end{table}

\end{document}
